% =====================================================================================
% Uporedna analiza reprezentacija pri simuliranim degradacijama slike
% primenjena na zadatak klasifikacije saobraćajnih znakova
%
% SKELET. Sadržaj se popunjava iz docs/report-material/*.md -- svaka sekcija ispod nosi
% komentar sa izvorom (nota + zadatak) i imenom slike iz figures/report/.
%
% Formatiranje je namerno identično predlogu projekta
% (docs/proposition/predlog_projekta.tex): isti documentclass, iste margine, isti stil
% sekcija (\section* sa ručnim brojem), obična bibtex bibliografija, isti blok za uklanjanje
% naslova "References". Bez babel paketa -- predlog ga ne koristi, a utf8 + T1 + lmodern je
% dovoljno za srpsku latinicu i izbegava zavisnost od serbian.ldf, koji na ovoj mašini nije
% instaliran (i koji po podrazumevanim podešavanjima daje ĆIRILICU).
% =====================================================================================

\documentclass[11pt,a4paper]{article}

\usepackage[utf8]{inputenc}
\usepackage[T1]{fontenc}
\usepackage{lmodern}
\usepackage{geometry}
\usepackage{enumitem}
\usepackage{hyperref}
\usepackage[dvipsnames]{xcolor}
\usepackage{array}
\usepackage{soul}
\usepackage{amsmath}
\usepackage{amssymb}
\usepackage{graphicx}
\usepackage{booktabs}
\usepackage{float}
\usepackage{caption}
\usepackage{subcaption}

\geometry{margin=2.2cm}
\hypersetup{
  colorlinks=true,
  linkcolor=black,
  urlcolor=blue,
  citecolor=black
}
\setlist[itemize]{leftmargin=1.2cm,itemsep=2pt,topsep=4pt}
\setlist[enumerate]{leftmargin=1.0cm,itemsep=4pt,topsep=6pt}
\setlength{\parindent}{0pt}
\setlength{\parskip}{6pt}
\renewcommand{\arraystretch}{1.2}

\renewcommand*\figurename{Slika}
\renewcommand*\tablename{Tabela}
\captionsetup[figure]{labelfont={bf},name={Slika},labelsep=colon}
\captionsetup[table]{labelfont={bf},name={Tabela},labelsep=colon}

\graphicspath{{../../figures/report/}}

% Skraćenice za imena metoda, da ostanu dosledna kroz ceo rad.
\newcommand{\pcasvm}{\texttt{pca\_svm}}
\newcommand{\hogsvm}{\texttt{hog\_svm}}
\newcommand{\bovwsvm}{\texttt{bovw\_svm}}
\newcommand{\cnnfeat}{\texttt{cnn\_feat\_svm}}
\newcommand{\cnne}{\texttt{cnn\_e2e}}

% Ukloni automatski naslov ("References") iz thebibliography okruženja -- isto kao u predlogu.
\makeatletter
\renewenvironment{thebibliography}[1]
     {\list{\@biblabel{\@arabic\c@enumiv}}%
           {\settowidth\labelwidth{\@biblabel{#1}}%
            \leftmargin\labelwidth
            \advance\leftmargin\labelsep
            \@openbib@code
            \usecounter{enumiv}%
            \let\p@enumiv\@empty
            \renewcommand\theenumiv{\@arabic\c@enumiv}}%
      \sloppy
      \clubpenalty4000
      \@clubpenalty \clubpenalty
      \widowpenalty4000%
      \sfcode`\.\@m}
     {\def\@noitemerr
       {\@latex@warning{Empty `thebibliography' environment}}%
      \endlist}
\makeatother

\begin{document}

\begin{center}
{\Large Univerzitet u Nišu\\Elektronski fakultet}\\[1.2cm]
{\LARGE \textbf{Projekat iz predmeta Računarski vid}}\\[0.6cm]
{\Large \textbf{Uporedna analiza reprezentacija pri simuliranim degradacijama slike primenjena na zadatak klasifikacije saobraćajnih znakova}}\\[1.2cm]
{\large Student: Mihajlo Madić, broj indeksa 2119}\\[0.3cm]
{\large Niš, 2026.}
\end{center}

\vspace{0.6cm}

% =====================================================================================
\section*{Sažetak}
% Izvor: README.md + PROJECT_TASKS.md §1 i §2. ~200 reči.
% Mora da sadrži: (1) GTSRB je zasićen benchmark, pa pitanje nije "koja reprezentacija je
% najtačnija" nego "koja otkazuje pod kojim uslovima"; (2) četiri reprezentacije, fiksiran
% linearni klasifikator, pet konfiguracija; (3) tri degradacije x pet nivoa; (4) glavni
% nalaz -- popuniti iz zadatka 9.8.

\noindent\textbf{Ključne reči:} računarski vid, klasifikacija saobraćajnih znakova, GTSRB,
reprezentacije slike, PCA, HOG, vreća vizuelnih reči, konvolucione neuronske mreže, robusnost.

% =====================================================================================
\section*{1. Uvod}
% Izvor: predlog §4.1-4.2, README.md.

\subsection*{1.1 Kontekst -- mesto modula u širem sistemu}
% Izvor: predlog §4.1. Ceo lanac: detekcija -> klasifikacija -> praćenje -> procena
% udaljenosti. Eksplicitno označiti koji JEDAN modul ovaj rad isporučuje.

\subsection*{1.2 Formulacija problema}
% Izvor: predlog §4.2. Ulaz: RGB isečak 15x15 do 250x250 sa tačno jednim znakom.
% Izlaz: jedna od 43 klase.

\subsection*{1.3 Istraživačko pitanje i doprinos}
% Izvor: PROJECT_TASKS §1. Doprinos NIJE rang-lista nego INTERAKCIJA reprezentacija x
% degradacija. Navesti falsifikabilnu tvrdnju iz §2 (ukrštanje PCA/HOG pod šumom i gamom).

\subsection*{1.4 Struktura rada}

% =====================================================================================
\section*{2. Reprezentacije slike}
% Izvor: note 11, 12, 13, 14 -- svaka ima sekciju "What the representation is".

\subsection*{2.1 Strukturne ose poređenja}
% Izvor: PROJECT_TASKS §1 tabela. OKOSNICA CELOG RADA -- tabela sa 4 reda (prostorna
% struktura / lokalnost / da li uči) mora doći rano, jer sva predviđanja i cela diskusija
% referišu na nju.

\subsection*{2.2 Analiza glavnih komponenti (PCA)}
% Izvor: note 11 §1-§2. Eigenfaces~\cite{eigenfaces}. Holistička, kritična na poravnanje.
% Odluka whiten=False i obrazloženje (note 11 §2.1).

\subsection*{2.3 Histogrami orijentisanih gradijenata (HOG)}
% Izvor: note 12 §1-§2. \cite{hog}. Kruta rešetka; blokovska L2-Hys normalizacija je
% mehanizam na kome počiva predviđanje o gami.

\subsection*{2.4 Vreća vizuelnih reči (BoVW)}
% Izvor: note 14 §1-§2. \cite{bovw,sift}. JEDINA reprezentacija bez prostornog rasporeda.
% Gusti SIFT umesto detektora i zašto (detektor vraća promenljiv broj tačaka po slici).

\subsection*{2.5 Konvoluciona neuronska mreža (CNN)}
% Izvor: note 13 §1-§2. \cite{ciresan}. Jedna mreža -> dve metode (\cnne{} i \cnnfeat{}).
% Slika: cnn_architecture.png. Zašto \cnnfeat{} uopšte postoji (note 13 §1).

% =====================================================================================
\section*{3. Skup podataka}

\subsection*{3.1 GTSRB}
% \cite{gtsrb}. 43 klase, 39.209 / 12.630 isečaka.

\subsection*{3.2 Struktura traka i podela na skupove}
% Izvor: note 04 -- NAJVAŽNIJA METODOLOŠKA SEKCIJA.
% Ključ je (class_id, track_id), NE sirovi TTTTT prefiks (75 naspram 1307).
% Izmereno: korelacija 0,61 unutar trake naspram 0,17 između traka iste klase.
% Slike: split_leakage.png, split_similarity.png.

\subsection*{3.3 Neuravnoteženost klasa}
% Izvor: note 02. 10,7x (2250 naspram 210) -- razlog zašto se izveštava makro-F1.

% =====================================================================================
\section*{4. Metodologija}

\subsection*{4.1 Pretprocesiranje}
% Izvor: note 08. Tri konfiguracije kao lestvica; svaka dodaje tačno jedan faktor.
% Slika: preprocessing_configs.png.

\subsection*{4.2 Kontrolisane degradacije}
% Izvor: note 10. Tri stresora x pet nivoa.
% Slike: degradation_contact_sheet.png, pipeline.png.
% KLJUČNO: degradacija se primenjuje POSLE pretprocesiranja -- ubacuje se rezidual koji
% pretprocesiranje nije uklonilo. Obrazloženje izbora baš ova tri stresora.

\subsection*{4.3 Fiksiran klasifikator i protokol izbora hiperparametara}
% Izvor: note 11 §6.2 i §7.4 (pitanje Q4).
% "Fiksiran klasifikator" = isti estimator i isti protokol izbora, NE isti hiperparametri.
% C i class_weight se biraju po metodi, na validaciji, po makro-F1.

\subsection*{4.4 Metrike}
% Izvor: note 05. labels=range(43). Makro-F1 je i kriterijum IZBORA, ne samo izveštavanja.

\subsection*{4.5 Merenje cene}
% Izvor: note 06. DVA broja za inferenciju (paketna i pojedinačna): PCA 49x, HOG 1,3x,
% CNN 0,87x. Veličina modela ne meri istu stvar za sve metode (PCA je 96% baza, HOG nema
% naučenu komponentu). Vremena su relativne cene u jednom okruženju.

\subsection*{4.6 Reproducibilnost}
% Izvor: note 03, note 07. Jedan seed, pinovanje niti, sadržajno izvedeni seedovi za
% degradacije, results.csv sa run_id. Granice tvrdnje (float32, raspored u memoriji, EPP).

% =====================================================================================
\section*{5. Implementacija}
% Izvor: note 01, 07, 08. Onoliko koliko treba da rezultat bude ponovljiv, ne više.

\subsection*{5.1 Organizacija koda i tok izvršavanja}
% Slika: pipeline.png.

\subsection*{5.2 Keširanje pretprocesiranih slika}
% Izvor: note 08. 83x ubrzanje; keš se proverava bit-po-bit prema referentnoj implementaciji.

% =====================================================================================
\section*{6. Eksperimentalna postavka}
% Izvor: PROJECT_TASKS §7. Mreža 5 metoda x 16 uslova = 80 evaluacija.
% Obuka je uvek na čistim podacima -- bez augmentacije i bez degradacija.

\subsection*{6.1 Predviđanja zabeležena pre eksperimenata}
% Izvor: predictions.md, datirano 13.09.2026, pre bilo kakvog modela.
% Tabela ide OVDE, a ne u diskusiju -- poenta je da su zapisana unapred.

% =====================================================================================
\section*{7. Rezultati}

\subsection*{7.1 Uporedna tabela metoda}
% Izvor: zadatak 9.1, results.csv. Pet redova.
% UPOZORENJE (note 06): kolona "veličina modela" ne meri istu stvar za sve metode.

\subsection*{7.2 Matrice konfuzije}
% Izvor: zadatak 9.2. Slika: confusion_best_worst.png.
% NAPOMENA: prikazati VIŠE od najbolje i najgore metode -- PCA i CNN dele iste najteže
% klase, a HOG ima potpuno druge (note 12 §5.2).

\subsection*{7.3 Najčešće zamene klasa}
% Izvor: zadatak 9.3, note 12 §5.2, note 13 §7.3.
% PCA i CNN: znaci prestanka zabrane (60,0% i 15,0%). HOG: ograničenja brzine (80->50, 16,9%).

\subsection*{7.4 Tačnost u funkciji veličine znaka}
% Izvor: zadatak 9.4. Slika: accuracy_by_size.png. 47,6% slika je ispod 32 px.

\subsection*{7.5 Krive robusnosti}
% Izvor: zadatak 9.5. Slika: robustness_curves.png -- GLAVNA SLIKA RADA.
% Relativno u odnosu na sopstvenu čistu osnovu svake metode.

\subsection*{7.6 Ablacija pretprocesiranja}
% Izvor: zadatak 9.7, note 08 dodatak. Izveštava se kao STABILNOST RANGIRANJA.
% Mora nositi nalaz o obavijanju nijanse (hue wrap) -- videti 8.5.

% =====================================================================================
\section*{8. Diskusija}

\subsection*{8.1 Predviđanja naspram ishoda}
% Izvor: zadatak 9.6 -- JEZGRO DISKUSIJE. Red koji kaže "pogrešno, jer..." vredi više.

\subsection*{8.2 Zašto reprezentacije otkazuju onako kako otkazuju}
% Izvor: note 11 §3.1 (PC1 je osvetljenje, 52,3%), §9.2 (PC2 je oblik), note 12 §3.1-3.2
% (linearni kontrast se poništava TAČNO, gama samo delimično).
% Slike: pca_eigensigns.png, pca_reconstruction.png, pca_pc1_brightness.png,
% hog_visualization.png.

\subsection*{8.3 Curenje podataka kroz trake -- izmereno, ne tvrđeno}
% Izvor: note 04, zadatak 4.5. Slika: split_leakage_measured.png.
% +5,08 p.p. tačnosti i +9,58 p.p. makro-F1. Makro-F1 se naduvava DVOSTRUKO.

\subsection*{8.4 Objašnjena varijansa nije diskriminativnost}
% Izvor: note 11 §3.2, §7.1c, §7.3. Slika: pca_component_sweep.png. Tri nezavisna slučaja.

\subsection*{8.5 Defekt u konfiguraciji clahe\_hsv -- nijansa se obavija}
% Izvor: note 08 dodatak. Crvena boja je na tački obavijanja (0/179): Stop ima 56,5%
% zasićenih piksela pri nijansi <10 i 4,4% pri >169. 3,98% susednih gradijenata nijanse
% prelazi 90. Cena raste sa oslanjanjem na izvode: HOG -9,7 p.p., CNN -2,3 p.p. (i vrhunac
% u epohi 5 umesto 18), PCA -3,0 p.p. NE sme se izvestiti kao "boja ne pomaže".

\subsection*{8.6 Ograničenja}
% Izvor: note 03 (jedan seed), note 06 (nema GPU; EPP menja vremena za 20%), note 09
% (zašto jitter nije moguć na GTSRB), note 11 §2.1 (nebeljena PCA), note 13 §1.1 (obeležja
% CNN-a optimizovana za softmax glavu), note 13 §2.2 (glava mreže manja zbog uporedivosti),
% note 14 §6.3 (BoVW bez prostorne piramide -- namerno).

% =====================================================================================
\section*{9. Zaključak}
% Izvor: zadatak 9.8 (pet konkretnih nalaza). Odgovoriti EKSPLICITNO na pitanje iz §1.3:
% da li se rangiranje menja između stresora. Oba ishoda su rezultat.

\subsection*{9.1 Pravci daljeg rada}
% Izvor: PROJECT_TASKS §11 (GTSDB, jitter) i §12 (degradacija u trenutku snimanja,
% složeni scenariji). Plus: ispravka clahe_hsv (nijansa kao cos/sin ili CIELab).

% =====================================================================================
\section*{10. Literatura}

\bibliographystyle{unsrt}
\nocite{szeliski, milosavljevic, laganiere}
\bibliography{references}

\newpage

\subsection*{Dodatak A -- Tabela akronima}
% Preneti iz predloga (Dodatak B) i dopuniti: GTSDB, CLAHE, SIFT, BoVW, HOG, PCA, SVM,
% CNN, ROI, BLAS, SMT, EPP.

\subsection*{Dodatak B -- Okruženje i verzije biblioteka}
% Izvor: note 01 + results/platform_<run_id>.json. Tačne verzije, broj niti, commit.

\end{document}
